\section{Discussion}
\label{sec:discussion}

\paragraph{Potential applications.}
Our findings suggest a way to evaluate the reuse of contextual skills.
Consider a Web agent that repeatedly follows the same instructions across
requests. A reusable internal representation would need to preserve task
performance in the receiving execution. Our results motivate evaluating
this transfer separately from the original execution's dependence on the
skill. If a reuse method also removes access to the original instructions,
that additional change needs its own evaluation. The observed cache gains
motivate reuse methods that carry both hidden states and attention caches
across requests. Deployment studies can evaluate their latency, memory
use and task performance in interactive Web settings.

\paragraph{Future work.}
A next step is to investigate which questions benefit from cache
replacement. The mathematical analysis connects remaining cache
differences to output stability under explicit assumptions. Measuring
these differences, their influence on token predictions and the donor's
prediction margins could test whether this framework helps distinguish
successful from unsuccessful repairs. The complete-source controls further motivate testing which parts of
the receiving history must be recomputed, and whether the required budget
can be predicted on held-out skills. Extending these controls beyond
Qwen3-8B/MedCalc would test their cross-model and cross-task scope. Separating attention cuts during prompt
processing from those during answer generation would clarify when continued
skill access matters. Finally,
multi-step Web workflows would assess how far the observed separation
extends beyond the question-answering settings studied here.
